% Timaeus (timaeus) --- English translation
% Translated by Claude (Anthropic) from the Latin (Perseus).
% Style guide: STYLE.md.

\ciceroSection{1} I have written much in my \textit{Academica} against the natural philosophers, and have often debated with Publius Nigidius in the manner and method of Carneades. For that man was furnished with all the other arts that befit a free man, and was besides a keen and careful investigator of those things that nature seems to keep hidden away; and in the end my judgment is this: that after those famous Pythagoreans, whose teaching had in a sense died out, though it had flourished for several centuries in Italy and Sicily, it was he who arose to bring it back to life.

\ciceroSection{2} When he was waiting for me at Ephesus as I was setting out for Cilicia — he himself returning to Rome from his post — and when Cratippus too had come to the same place from Mytilene to greet me and to see me — Cratippus, in my judgment easily the first of all the Peripatetics I have heard — I saw Nigidius with great pleasure and came to know Cratippus. And the first stretch of time, the time of greeting, we spent in inquiry . . .

\ciceroSection{3} What is it that always is and has no coming into being, and what is it that comes into being and never is? The one is grasped by intelligence and reason, being always one and the same; the other, which presents itself to opinion by way of sense without reason, and is wholly a matter of opinion, comes into being and perishes and can never truly be. But everything that comes into being must of necessity come into being from some cause; for the origin of nothing can be found if its cause is removed.

\ciceroSection{4} Therefore, if the one who undertakes to make some work looks to the pattern that is always the same and sets that before himself as his model, he must of necessity produce a splendid work; but if he looks to the pattern that comes into being, he will never attain the beauty he is reaching for. The whole heavens, then, or the world — or whatever other name it delights to bear, by this let it be called by us —

\ciceroSection{5} concerning it let us first consider the question that must be considered at the outset of every inquiry: whether it always was, generated by no coming into being, or whether it has come into being from some beginning in time. It has come into being, since it is seen and touched and is on every side endowed with body. But all such things move the sense, and the things that move the sense are the same that settle in opinion; and these, we have said, have a coming into being and are generated, and nothing can come into being without causes.

\ciceroSection{6} And indeed to find that one who is, as it were, the parent of this whole is difficult, and, once you have found him, to declare him to the crowd is forbidden. Again, then, it must be considered whether that maker of so great a work imitated the pattern that is always one and the same and like itself, or the one that we call generated and come into being. And surely, if this world is beautiful and its craftsman good, beyond doubt he chose rather to imitate the pattern of eternity; but if otherwise — which it is not even permitted to say — he followed a generated model in place of an eternal one.

\ciceroSection{7} There is no doubt, then, that he chose rather to pursue eternity, since nothing is more beautiful than the world and nothing more excellent than its builder. So, then, it was generated, fashioned after that which is grasped by reason and wisdom and held fast by unchanging eternity. From which it follows that this world we behold must of necessity be the eternal image of something eternal. But the most difficult thing in every search of reason is the beginning. Of the things we have said, then, let this be the first distinction.

\ciceroSection{8} Every discourse, it seems, has a kinship with the things it sets forth. And so, when the discourse argues about what is stable and unchanging, let it be such as that subject is — such as can neither be refuted nor convicted. But when it has entered upon what is imitated, upon fashioned images, it thinks it does well if it attains a likeness of the truth. For as much as eternity is to that which has come into being, so much is truth to belief. And so, if perhaps in discoursing on the nature of the gods and the coming into being of the world we attain less than our heart desires — namely that the whole discourse should be clear and plainly set out, consistent with itself and in every part in agreement with itself — it will surely be no wonder, and you ought to be content if probable things are said. For it is fair to remember both that I, who shall discourse, am a man, and that you, who judge, are men, so that, if probable things are said, you require nothing beyond them.

\ciceroSection{9} Let us ask, then, the cause that moved the one who contrived these things to seek the origin of all things and a new making. He excelled, plainly, in goodness; and the good man envies no one; and so he generated all things like himself. This, surely, was the most just cause of the world's coming into being. For when god had determined to fill the world with all good things and to mingle nothing of evil in it, so far as nature would allow, he took to himself whatever fell within the perception of sight — not tranquil and at rest, but stirred without measure and adrift — and he led it from disorder into order; for this he judged to be the better.

\ciceroSection{10} And it is not permitted, nor ever was, that the one who is best should make anything but the most beautiful. So when he had taken thought, he found that of the things perceived by nature nothing without intelligence could in its whole kind be more excellent than that which has intelligence. For which reason he enclosed intelligence within soul, and soul within body. Thus he reckoned that work to be made most beautiful. For this cause it must not be hesitated to declare, if only something can be traced out by conjecture, that this world is a living being, and that intelligent and established by divine providence.

\ciceroSection{11} This being laid down, what follows must be considered: in fashioning the world, the likeness of which living being god followed. Of none, surely, of those living beings known to us; for they are all divided into certain kinds, or are merely begun and in no part complete; but nothing can be beautiful that resembles what is imperfect and unfinished. Let us say, then, that the world resembles that of which every living being is, as it were, a small particle, whether it be discerned in single instances or in the whole kind. For all the living beings that are discerned by the mind and grasped by reason are embraced within the compass of reason and intelligence, just as men and cattle and all things that fall under sight are embraced by this world.

\ciceroSection{12} For that which can be conceived as most beautiful in the nature of things, and which is in every part most complete — when god wished to make the world its likeness, he made one visible living being, in which all living beings might be contained. Have we then rightly said that the world is one, or was it better and truer to say there were several or countless worlds? One, surely, if indeed it was made after the pattern. For that which embraces all living beings that are grasped by reason cannot be together with another. For again there would have to be another living being to contain it, of which the former living beings would be parts, and these heavens would be the image of that last one, not of the nearer. So that none of this might come about, and that this world might be most like the complete living being, for this very reason — that it was alone and one — god therefore brought forth this world single and only-begotten.

\ciceroSection{13} But everything that has come into being must of necessity be bodily and visible and likewise tangible. Further, nothing can be looked upon and seen that is empty of fire, nor indeed touched that lacks the solid, and nothing solid that is without earth; for which reason god, in laboring to make the world, first joined earth and fire. But all two things require some third thing for their cohesion, and seek, as it were, a knot and bond. And of bonds that one is most fitting and most beautiful which makes itself and the things it binds together as much one as possible. This is best achieved by what in Greek is called \textit{[Greek: analogia]}, and which in Latin — for one must be bold, since these things are first being coined anew by us — may be called comparison or proportion.

\ciceroSection{14} For whenever it happens, among three numbers or figures or things of any kind whatever, that the one which is the middle stands to the first in the same proportion in which the last stands to it, and again that, as the last is to the middle, so the middle is to the first, then the term that is middle becomes at once first and last, while the last and the first become middle terms; in this way necessity compels that those which had been distinct should be the same, and once they have been made the same, the result is that all are one. And if the body of the whole had been laid out flat and even, with nothing . . . in it, a single middle term set between would have bound together both itself and the things between which it was placed.

\ciceroSection{15} But since solidity was sought for the world, and since solid things are never joined by one middle term but always by two, it so came about that god set water and air between fire and earth, and brought them into proportion with one another and joined them by analogy, so that, as fire is to air, so air is to water, and as air is to water, so by the same proportion water is to earth. From this conjunction the heavens have been so fitted together that they fall under sight and touch. And so for this reason, and out of these elements four in number, the body of the world was wrought, bound by the proportion I have described; whence it embraces itself in a kind of concordant friendship and affection, and so aptly coheres that it can in no way be dissolved except by the one by whom it was bound together.

\ciceroSection{16} And of those four things which I named above, all the parts have been so disposed throughout the whole world that no part of any one of these kinds was left outside, and that within this whole all those kinds were present entire; and this for these reasons: first, that the world, being a living creature, might be made perfect out of perfect parts; next, that it might be one, with no part left over from which another might be generated; and finally, that no disease and no old age might be able to touch it.

\ciceroSection{17} For every framing-together of a body is shaken and broken by the force of heat or of cold or by some violent impact, and is driven on toward diseases and old age. This, then, was the reckoning that god, the maker and builder of the world, kept in view: that he should bring to completion a single work, whole and perfect, out of elements all whole and perfect, a work free from every disease and from age. And he gave it a shape both most akin to it and most fitting. For from the living creature by which he willed all the rest of living creatures to be contained, he fashioned this one in that shape within which alone all other shapes are enclosed, and he made it spherical, which the Greeks call \textit{[Greek: sphairoeides]}, every point of whose surface is reached by equal radii from the center; and he turned it on the lathe in such a way that he could make nothing rounder, that it might have nothing of roughness, nothing of obstruction, nothing cut into angles, nothing of windings, nothing jutting out, nothing hollowed, and all its parts most alike to all the others — for to his judgment likeness was better than unlikeness.

\ciceroSection{18} And he surrounded the whole figure of the world all over with smoothness. For it had no need of eyes, since nothing was left outside that could be seen, nor of ears, since there was nothing even to be heard, nor was there breath poured around the outer edges of the world, that it should require respiration; nor indeed did it want either nourishment of the body or the removal of food digested and consumed; for no departure could take place and no addition, nor was there anywhere for these to come from. And so it fed itself by the wasting and aging of itself, since by itself and through itself and from itself it both suffered and did all things. For so he reckoned, the one who joined and founded these things, that the world should be content with itself and have no need of anything else.

\ciceroSection{19} And so he attached to it no hands, since there was nothing to be grasped and nothing to be repelled, nor feet nor other limbs by which it might support the body's walking. For he gave to the heavens the motion that was most fitting to its shape, the one which of the seven motions stirs mind and intelligence most. And so by one and the same revolution it is turned and spun about itself. The remaining six motions he kept apart from it, and thus freed it from all wandering. For this revolution, then, which had no need of feet or stepping, he gave it no limbs for walking.

\ciceroSection{20} This god who already was, thinking of a god that would one day be, made it smooth and even on every side, equal from the center to the edge, and perfect and complete out of things complete and perfect. And as he set the soul in the middle of it, so he stretched it through the whole; then he surrounded it with body and clothed it from without, and embraced it with a heaven solitary in its wandering and revolving and driven round in a circle, one which by reason of its excellence could readily keep company with itself and want no other, being sufficiently known and familiar to itself.

\ciceroSection{21} Thus that eternal god brought forth this god, perfectly blessed. But he did not begin the soul, as we just now spoke of it, only then, when he had made the body for it; for it would not be right that the greater should obey the lesser. But we say many things rashly and without consideration. God, however, begot the soul both prior in origin and prior in excellence, and set it over the body as its master and ruler over what obeys; and he wrought this in some such manner as this. Out of that matter which is indivisible and always of one mode and like itself, and out of that which is generated as divisible in bodies, he mingled into the middle a third kind of matter from the two, which should be both of the same nature and of the other; and this he set between the indivisible and that which was divisible in body.

\ciceroSection{22} When he had taken these three, he tempered them into a single form, and that nature which we have called the other he joined by force with the same, though it shrank away and was averse to that coupling; and mingling these with matter, when he had made one out of three, he divided that very thing into the parts that were fitting. Now he tempered the single parts out of the same and out of the other and out of matter. And this was the division: first he drew off one part from the whole; then a second, double the first; then a third, which was one and a half times the second and triple the first; then a fourth, which was double the second; then a fifth, which was triple the third; then a sixth, eight times the first; and last a seventh, which exceeded the first by twenty-seven parts.

\ciceroSection{23} Then he set about filling out the double and triple intervals, cutting parts again from the whole; and these he so placed within the intervals that in each there were two means (for I scarcely dare say "medieties," which the Greeks call \textit{[Greek: mesotētas]}; but let it be understood as though I had so said, for it will be plainer that way): one of them exceeding the extremes by the same part by which it was exceeded, the other exceeding the extremes by an equal number and exceeded by an equal number. And taking from these conjunctions intervals of one and a half, of one and a third, and of one and an eighth, he filled out all the intervals of one and a third within the first intervals by an interval of one and an eighth, leaving over a little portion of each.

\ciceroSection{24} And with the interval of that little portion left over, number had to number, among the extremes, the same proportion and ratio that two hundred and fifty-six have to two hundred and forty-three; and so that mixture out of which he cut these parts he had now wholly used up. This whole conjunction, then, he split lengthwise into two, and fitting middle to middle he crossed them as in the letter X; then he bent them into a circle, so that they should be joined both each to itself and to each other at the junction lying opposite, and by that motion whose circle was always in the same place and stirred in the same way, he embraced them on every side.

\ciceroSection{25} And so, when he had set the one to embrace the outer circle and the other the inner, he marked the former as of the nature of the same, the latter as of the other; and that which was of the same he turned aside from the flank toward the right, while this nearer one he directed from the median line toward the left, but he gave sovereignty to the upper, which alone he left undivided. The inner he divided into six parts, and bade seven unequal circles be moved at intervals of double and triple in courses contrary to one another. Of these he made three equal in speed, but four both unequal among themselves and unlike the remaining three.

\ciceroSection{26} When, then, that god who brought forth the world had begotten the soul out of his own mind and will, only then did he spread beneath the soul, and make inner to it, all that was concrete and corporeal, and so, fitting middle to middle, he coupled them together. Thus the soul, setting out from the center, threw itself in a round circuit about the outermost edge of the heavens from the highest region, and, turning upon itself, ushered in the divine beginning of an everlasting and wise life.

\ciceroSection{27} And the body of the heavens, indeed, was made visible, but the soul escapes the gaze of the eyes. It is, of all things, the one that shares in and partakes of reason and of that concord which the Greeks call \textit{[Greek: harmonia]}, of the eternal things that fall under intelligence; than which nothing better has been generated by the best and most excellent begetter. For, joined out of the same nature and the other, with matter set alongside and compacted by the proportionate tempering of three parts, it turns itself about; and when it has laid hold on changeable matter, and again on what is indivisible and simple, it is moved through itself entire and discerns what is of the same kind and what is of the other, and judges all the rest — what is most fitting to each thing, what falls to each in its place or manner or time, and what distinction there is between the things that come into being and the things that are forever the same.

\ciceroSection{28} But true reason, which is concerned with the things that are forever the same and with the things that change, when it is moved within the same and within the other, by itself, without voice and without any sound, when it touches that part by which sense can be stirred, and the circuit of that other kind, unaltered and straight, declares all things to the soul and the mind, then there are generated firm and true opinions and assents; but when it turns among those things which, abiding ever the same, are held together not by sense but by intelligence . . . By reason, then, and by the divine mind, for the origin of time the course of the sun and the moon was devised

\ciceroSection{29} . . . nature would turn, so that the course of the moon would circle nearest the earth, and next above the earth would be the revolution of the sun. Then the Morning Star and the holy star of Mercury hold a course equal in swiftness to the sun, but a certain force contrary to it; and by that running together which the Morning Star, Mercury, and the sun have among themselves, the one outstrips the others and is outstripped by them in turn.

\ciceroSection{30} What the reason was for placing the remaining heavenly bodies, and what their placement is, must be deferred to another discourse, lest a longer account be set down upon the matter that had to be touched on than upon the matter for whose sake we touch on it. When, therefore, each of those heavenly bodies had obtained the seemly course from which the measure of time was to be marked out, and when, their bodies being bound together by living bonds, living beings came into being and learned to obey command, then, running across by the motion of the other nature into the motion of the same nature, and clinging there and held fast within it, since some traversed a larger orbit and some a smaller — those a larger more slowly, those a smaller more swiftly — by the motion of the one and the same nature the ones that moved most swiftly seemed to be outstripped in speed by the slower, and, when they were overtaking, to be overtaken.

\ciceroSection{31} For it turned every orbit of theirs by a kind of spiral inflexion, because the two movements proceeding contrariwise at once brought it about that what was slowest came nearest to the swiftest. And so that there might be some clear measure to declare the swiftnesses and slownesses in the eight courses, the god himself kindled the sun like a lamp at the second circuit above the earth, so that the heavens might shine out as brightly as possible upon all things, and that the living beings which had the right to be taught might learn the nature and force of numbers from the motion of the same and of that which was like it.

\ciceroSection{32} Night, then, and day, generated in this manner and for these causes, make up one circuit of the most wise and best orbit; and the month, when the moon, having traversed her course, has caught up with the sun; and the year, when the sun has completed and traveled over his whole orbit.

\ciceroSection{33} But the circuits of the other heavenly bodies, of which men are ignorant save for very few, they neither call by name nor measure off against one another by number. And so they do not know that these wanderings of the heavenly bodies are that very thing which is rightly called time, endowed as they are with infinite multitude and wondrous variety. And yet this much can be perceived and understood: that, by the absolute and perfect number of time, the absolute and perfect year is then at last completed, when the eight circuits, their courses being run, have returned to the same starting point, and when one and the same orbit, ever like itself, has measured them all out.

\ciceroSection{34} For these causes, then, were the stars born which, passing through the heavens, should turn themselves back by their solstitial and wintry return, so that this whole living being which we see might be most like to that living being which we apprehend by thought, in imitation of eternity. And the rest, indeed, up to the origin of time, he had molded as imprinted from those things which it imitated; but because he had not yet enclosed every living being within the world, in that respect the likeness of the image fell short of the pattern set before it. As many, then, and of such kinds, as were the forms of living beings that the mind, gazing upon the species of things, could discern, so many and of such kinds did he resolve to fashion within himself in this world.

\ciceroSection{35} Now there were four kinds of living beings: one divine and heavenly, the second winged and of the air, the third aquatic, the fourth of the earth. The form of the divine kind of life he made chiefly of fire, that it might be most splendid and most beautiful to look upon. And since he wished to make it like the nature of the whole, he rounded it into a sphere for the sake of its turning, and made it the companion of the wisdom of the best mind, and distributed it equally about the whole heavens, so that this thing, marked out with variety, the Greeks well call \textit{[Greek: kosmon]}, and we the shining world.

\ciceroSection{36} To the divine kinds he gave two kinds of motion: one, that it should always be in the same place . . . and revolve about the same axis for all things and in one manner; the other, that it should be driven forward by the revolution of the same and the like. But of the five remaining motions he willed it to be bereft and without share, unmoving and standing still. Of this kind are those heavenly bodies which, fixed in the heavens, do not move from their place, which are living beings, and divine, and for that cause cling to their seats and abide forever. But those which glide along in wandering and changeable straying were generated as we said above.

\ciceroSection{37} And the earth, indeed, our nurse, which is upheld by the axis driven through it, the maker of day and night and likewise their guardian, he willed to be the most ancient of all the gods that should be brought to birth within the heavens. But the sportings of the gods, and the encounterings among the gods themselves, and the revolutions and the advances that come about in their orbits, and how they almost touch one another — which of them are joined close, and in which opposite region and behind which others, or before which, they glide, and at what seasons they vanish from our sight and, emerging again, strike terror into those bereft of reason — if we should try to set these out in words with no image of those things placed before the eyes, the labor would be undertaken in vain. But let what we have said be enough; and let the things we have prefaced concerning the nature of the gods that are seen and that have come into being have this for their limit.

\ciceroSection{38} But concerning the rest, whom the Greeks call \textit{[Greek: daimonas]}, and ours, I suppose, the Lares — if only this can rightly seem to have been so rendered — both to know and to declare their origin is a greater thing than that we should dare to profess to write of it. We must trust, no doubt, the men of old and of former times, as they say, who declared themselves the offspring of the gods, and so handed down to us their names. They could surely know their own begetters best of all, and it is hard not to put faith in men sprung from gods, although their account is confirmed neither by proofs nor by sure reasonings; but since they seem to speak of things known to them, we must obey the ancient law and custom.

\ciceroSection{39} So then, as it has been handed down by them, let the origin of these gods be held and told: that we record Oceanus and Salacia as generated and brought forth by the sowing of Caelus and the conceiving of Terra; from these, Saturn and Ops; thereafter Jupiter and Juno, and the rest, whom we see treated and called brothers and kindred among themselves, and their lineage, to use the old word.

\ciceroSection{40} When, therefore, all the gods had been created—both those who move and openly show themselves, and those who are revealed to us only so far as they themselves will it—then the god who had begotten all things speaks to them thus: "Give heed, you who are sprung from the seed of the gods. Of these works I am the parent and the maker; and they cannot be dissolved against my will—though everything that has been bound together can be dissolved. Yet it is in no way the part of a good being to wish to undo what reason has bound. And since you have come into being, you cannot indeed be immortal and indissoluble; yet you shall by no means be dissolved, nor shall any fated death destroy you, nor shall any treachery prove stronger than my purpose, which is a greater bond for your perpetuity than those by which you were bound together at the time you were begotten.

\ciceroSection{41} Learn, then, what is in my mind. Three kinds remain for us to make, and these are mortal; and if they are passed over, the completion of the heavens will not be perfect. For it will not embrace all the kinds of living things; yet it must embrace them, so that nothing be wanting from the whole. If these were made by me myself, they could match the life of the gods. That they may be generated, then, under a mortal condition, do you undertake the task: bring them into being, and imitate that power of mine which you remember I used in your own coming-to-be. And among them, those who shall be created of such a kind that they ought to be, as it were, kinsmen of the immortal gods, let them be called of the divine race, and let them hold the rule over all living things and willingly obey you by right and by law; the beginning and the seed of these shall be handed over to you by me, and you for your part shall weave on, to that which is immortal, a mortal part. So shall living creatures arise, whom you shall both feed while they live and receive into your bosom when they are spent."

\ciceroSection{42} Thus he spoke. Then he returned to the earlier mixing, in which, tempering the whole soul of universal nature, he blended it together; and pouring out the remnants of the earlier mixture, he made them even in much the same way as before, except that they were not so uncorrupted as those which are always the same, but second and even third in degree from these.

\ciceroSection{43} When, therefore, the whole had been wholly established, he distributed a number of souls equal to the stars, and joined each to each, and so set them, as it were, upon the chariot of the universe; and he showed them the laws of fate and necessity, and revealed that there would be one first coming-to-be, the same for all, measured and constant and diminished by none; and that, the souls having been sown and scattered, as it were, there would arise at fixed intervals of time a living creature that should be most fitted for the worship of the gods.

\ciceroSection{44} But since the nature of the human race was twofold, the matter stood thus: that the more excellent kind was that of those who were destined to be men. And when he had implanted the souls in bodies by necessity, and there came to those bodies now an accession, now a withdrawal, it was first of all necessary that there should arise a single sense, common to all, roused by a more violent motion and bound up with their nature; then love, mingled of pleasure and pain; and after these, anger and fear and the other motions of the soul, the companions of the former and others set against them in discord. Whoever rules these by reason will live justly; whoever surrenders himself to them, unjustly.

\ciceroSection{45} And he who has rightly and honorably finished the course of living given him by nature shall return to the star with which he was paired; but he who has lived without measure and without restraint—a second birth shall transfer him into the form of a woman, and if not even then he makes an end of his vices, he shall be tossed still more harshly, and shall be transferred, after the likeness of his own character, into the forms of cattle and wild beasts; nor shall he see the end of his evils before he has begun to follow that revolution of the same and the like which was inborn and implanted within himself. And this shall come to pass when, by reason, he has driven off those turbulent and reason-bereft elements that have settled upon him out of fire, soul, water, and earth, and has come back to the first and best disposition of the soul.

\ciceroSection{46} When he had ordained these things in this way, and had set himself beyond all blame and cause of blame, if afterward any wrongdoing or fault should arise, he sowed them, as it were, scattering some into the earth, some into the moon, some into the other parts of the world, which are set as the signs and marks of the spaces of time. But after that sowing he committed it to the gods—younger ones, so to speak—to fashion mortal bodies, and to put the final polish and completion upon all of the human soul that remained, that ought to be added, and upon all that followed from it; and then to offer themselves as the chiefs and guides of this living creature, and to rule and govern its life as beautifully as could be, so far as it should not, being itself well made, seek out some misery for itself through its own fault.

\ciceroSection{47} And he indeed, who had composed all things, abode constantly in his own state; but those who had been created by him, when they had come to know the order of their parent, followed it. And so, when they had received the immortal principle of a mortal living creature, imitating their begetter and maker, they borrowed from the world particles of fire and earth and water and soul, which they should give back again, and joined these among themselves—not with the same bonds by which they themselves were bound together, but with bonds of such a kind as could not be seen on account of their smallness; and with these, like close-set little wedges fused together, they made of all the parts a single body, and into it, as the divine soul's circuit flowed in and flowed out, they bound that soul fast.

\ciceroSection{48} And so those circuits, plunged into the stream, neither held nor were held, but with great force now bore things along, now were borne along themselves. Thus the whole living creature was indeed moved, but without measure and at random, so that it was carried by the six motions. For both forward and backward, and to the left and to the right, and upward and downward, now this way, now that . . .

\ciceroSection{49} But if it has settled upon a bright surface, then either the same image or sometimes an altered one is returned, when the fire of the eyes has merged and joined itself with that fire which is spread before the face. And the things on the left appear on the right, because the contrary parts of the eyes touch the contrary parts. But right answers to right and left to left in the reflection of the lights, when these do not cling together. This happens when the smoothness of mirrors has taken on a height from this side and from that, and so has thrust the right of the eyes into the left part and the left into the right. Faces are seen upside down too, by the throwing back of the lights, which, turning them, renders the lower parts that which are the upper.

\ciceroSection{50} And all these things are of that kind which assist the causes of things; these are the ministrations the god employs, when, so far as it can be done, he brings about the appearance of the best. But most people suppose that these are not the helpers of causes, but are themselves the causes of all things—the things that have the power of cooling and heating, of condensing and dissolving, but lack all intelligence and reason, which are found nowhere in nature except in the soul. But the soul escapes wholly the sense of the eyes; whereas fire, soul, water, and earth are bodies, and these are seen.

\ciceroSection{51} But he who declares himself a lover of intelligence and wisdom must seek out the first causes of an intelligent and wise nature, and then the secondary causes of those things which necessarily move others, while themselves being moved by others. Wherefore I judge that we must proceed thus: that we should ourselves speak of both kinds of causes, but separately—of those which, together with intelligence, are the makers of the most beautiful and the best of things, and of those which, devoid of foresight, bring about results inconstant and disordered.

\ciceroSection{52} As for the causes of the eyes, that they should have the power they now have, I think enough has more or less been said. But their greatest usefulness, bestowed upon the human race by the gift of the gods, may now in turn be set forth. For it is the eyes that have brought us the knowledge of the best of things. For this discourse of ours about the universe would never have been discovered, if neither the stars nor the sun nor the heavens could have fallen under the gaze of our eyes. But as it is, day and night, made known to the eyes, and then the revolutions of months and years, have contrived the nature and power of number, and have measured out the span of time, and have driven us to the inquiry into the whole of nature; and from these things we have gained philosophy, than which no greater good has ever been given, nor shall be given, to the race of mortals by the granting and the gift of the gods . . .

